\begin{description}
\item[(A)] \textbf{Transiting Exoplanet (15 points).} \ioaafig{2025_TH_Q09-s1.pdf}
Measure the distance between successive minima, then determine the X-axis scale from the 5 and 10 day marks. The key's measurements are 5 days for 2.6\,cm and 2.5\,cm between minima, giving
\[
P=2.5\times\frac5{2.6}\approx4.81\text{ days}.
\]\hfill(5 points)
Printing or zoom can change the scale, but should preserve the period. Answers within $\pm0.1$ days receive full credit.

An alternative is to use a longer baseline. From 5 to 20 days is 7.8\,cm. From the first to the penultimate minimum is 10\,cm, spanning four periods. Thus
\[
P=\frac{10}{4}\frac{15}{7.8}\approx4.81\text{ days}.
\]
Both methods receive full credit. Let $n\%$ be the minimum plotted value. Then
\[
100\frac{L_{\mathrm{transit}}-L_{\mathrm{average}}}{L_{\mathrm{average}}}=n.
\]
With star radius $R$ and planet radius $r$,
\[
L_{\mathrm{transit}}=\frac{L_{\mathrm{average}}}{\pi R^2}(\pi R^2-\pi r^2).
\]
Substitution yields
\[
\left(\frac rR\right)^2=-\frac n{100},\qquad\frac rR=\sqrt{\frac{|n|}{100}}.
\]\hfill(5 points)
On the Y-axis, 0.6\% spans 2.2\,cm. A minimum's depth varies from 2.8 to 3.2\,cm. The marking key accepts 2.8--3.1\,cm for full credit, or an average over all minima. Using 2.9\,cm,
\[
|n|=2.9\times\frac{0.6}{2.2}\approx0.79\%.
\]\hfill(3 points)
Therefore $r/R\approx0.09$. \hfill(2 points)
The extreme values give 0.087 and 0.093, both rounding to 0.09.
\item[(B)] \textbf{Eclipsing Binary Star System (25 points).} The X-axis scale is 20 days to 9.1\,cm and the measured period spans 1.1\,cm. Thus
\[
P=1.1\times\frac{20}{9.1}\approx2.42\text{ days}.
\]\hfill(6 points)
Accept periods within $\pm0.1$ days. For a cooler large star of radius $R$ and hotter small star of radius $r$,
\[
L_{\mathrm{average}}=L_R+L_r.
\]
At secondary minimum the smaller star transits the larger,
\[
L_{\mathrm{Secondary}}=L_r+\frac{L_R}{\pi R^2}(\pi R^2-\pi r^2).
\]
At primary minimum, the smaller star is completely eclipsed, so $L_{\mathrm{Primary}}=L_R$. \hfill(6 points)
The luminosity ratios satisfy
\begin{align*}
n_{\mathrm{Primary}}&=\frac{L_R}{L_R+L_r},\\
n_{\mathrm{Secondary}}&=\frac{L_r+L_R(\pi R^2-\pi r^2)/(\pi R^2)}{L_R+L_r}\\
&=1-\frac{L_R}{L_R+L_r}\frac{r^2}{R^2}=1-n_{\mathrm{Primary}}\frac{r^2}{R^2}.
\end{align*}
Hence
\[
\frac rR=\sqrt{\frac{1-n_{\mathrm{Secondary}}}{n_{\mathrm{Primary}}}}.
\]\hfill(5 points)
The Y-axis scale is 0.2 to 1.1\,cm. The primary depth spans 1.8\,cm, so
\[
n_{\mathrm{Primary}}=1-1.8\frac{0.2}{1.1}=1-0.327\approx0.67.
\]\hfill(3 points)
The secondary depth spans 0.9\,cm, so
\[
n_{\mathrm{Secondary}}=1-0.9\frac{0.2}{1.1}=1-0.16\approx0.84.
\]\hfill(3 points)
Therefore
\[
\frac rR=\sqrt{\frac{1-0.84}{0.67}}\approx0.49.
\]\hfill(2 points)
Values within $\pm0.02$ receive full credit.
\end{description}
